% GENERATED FILE. Edit canonical lessons, then run npm run generate:books.
% canonical-source-hash: fnv1a64:7ef265240c3000d5
% canonical-lessons: JA-C407-sanpai, JA-C407-butsuzou, JA-C407-daibutsu, JA-C407-hondou, JA-C407-ogamu, JA-R407-review-402-407

\chapter{Inside the Temple}
\label{ch:ja-inside-the-temple}
\hlchaptermodality{407}

\begin{chapteropening}
\textbf{By the end of this chapter:} \emph{I can visit a temple and describe what is inside.}

Complete the last lesson of chapter 407: I can visit a temple and describe what is inside.
\end{chapteropening}

\section[\texorpdfstring{\ja{さんぱい} (sanpai)}{さんぱい (sanpai)}]{\ja{さんぱい} (sanpai) — a visit to worship (at a shrine or temple)}

\label{lesson:JA-C407-sanpai}



\begin{quote}
Before the new one: say the Japanese for a Shinto priest, then the Japanese for to offer up.
\end{quote}

\subsection*{You'll want to know: \ja{さんぱい}}
\textbf{\ja{さんぱい}} — \emph{sanpai} — \textquotedblleft{}a visit to worship (at a shrine or temple)\textquotedblright{}.

\begin{quote}
\textbf{\ja{かぞくで} \ja{てらに} \ja{さんぱいに} \ja{いきました。}} — \emph{kazoku de tera ni sanpai ni ikimashita} — \textquotedblleft{}My family went to the temple to worship.\textquotedblright{}
\end{quote}

\begin{cousinweb}
\textbf{\ja{さんぱい}} (\emph{sanpai}) is Sino-Japanese (kango): two Chinese characters read with their on'yomi, readings borrowed from Chinese. They are \emph{san}, \textquotedblleft{}visit\textquotedblright{} (Middle Chinese \emph{tshom}) and \emph{hai}, \textquotedblleft{}worship\textquotedblright{} (\emph{peajH}), joined as \emph{sanpai} (KANJIDIC2; Middle Chinese in Baxter's transcription).
\end{cousinweb}

\begin{grammarlens}[title={What you've built}]
One more word for getting around.
\end{grammarlens}

\subsection*{Your turn}
\begin{itemize}
\raggedright
  \item \emph{Say it:} \emph{sanpai}
  \item \emph{Say it:} \emph{sanpai}, once more
  \item \emph{Say it:} \emph{kazoku de tera ni sanpai ni ikimashita} — My family went to the temple to worship.
  \item \emph{Recall:} say the Japanese for a Shinto priest, then the Japanese for to offer up, then say \emph{sanpai} again
\end{itemize}

\begin{tcolorbox}[breakable,colback=teal!4,colframe=teal!35!black,title={Before you move on}]
What does \ja{さんぱい} mean? (\textquotedblleft{}A visit to worship (at a shrine or temple)\textquotedblright{}.) Say it once more.
\end{tcolorbox}
\section[\texorpdfstring{\ja{ぶつぞう} (butsuzō)}{ぶつぞう (butsuzō)}]{\ja{ぶつぞう} (butsuzō) — a statue of the Buddha}

\label{lesson:JA-C407-butsuzou}



\begin{quote}
Before the new one: say the Japanese for to offer up, then the Japanese for a visit to worship.
\end{quote}

\subsection*{You'll want to know: \ja{ぶつぞう}}
\textbf{\ja{ぶつぞう}} — \emph{butsuzō} — \textquotedblleft{}a statue of the Buddha\textquotedblright{}.

\begin{quote}
\textbf{\ja{てらで} \ja{ふるい} \ja{ぶつぞうを} \ja{みました。}} — \emph{tera de furui butsuzō o mimashita} — \textquotedblleft{}I saw an old statue at the temple.\textquotedblright{}
\end{quote}

\begin{cousinweb}
\textbf{\ja{ぶつぞう}} (\emph{butsuzō}) is Sino-Japanese (kango): two Chinese characters read with their on'yomi, readings borrowed from Chinese. They are \emph{butsu}, \textquotedblleft{}Buddha\textquotedblright{} (Middle Chinese \emph{bjut}) and \emph{zō}, \textquotedblleft{}statue\textquotedblright{} (\emph{zjangX}) (KANJIDIC2; Middle Chinese in Baxter's transcription).
\end{cousinweb}

\begin{grammarlens}[title={What you've built}]
One more word for getting around.
\end{grammarlens}

\subsection*{Your turn}
\begin{itemize}
\raggedright
  \item \emph{Say it:} \emph{butsuzō}
  \item \emph{Say it:} \emph{butsuzō}, once more
  \item \emph{Say it:} \emph{tera de furui butsuzō o mimashita} — I saw an old statue at the temple.
  \item \emph{Recall:} say the Japanese for to offer up, then the Japanese for a visit to worship, then say \emph{butsuzō} again
\end{itemize}

\begin{tcolorbox}[breakable,colback=teal!4,colframe=teal!35!black,title={Before you move on}]
What does \ja{ぶつぞう} mean? (\textquotedblleft{}A statue of the Buddha\textquotedblright{}.) Say it once more.
\end{tcolorbox}
\section[\texorpdfstring{\ja{だいぶつ} (daibutsu)}{だいぶつ (daibutsu)}]{\ja{だいぶつ} (daibutsu) — a great statue of the Buddha}

\label{lesson:JA-C407-daibutsu}



\begin{quote}
Before the new one: say the Japanese for a visit to worship, then the Japanese for a statue of the Buddha.
\end{quote}

\subsection*{You'll want to know: \ja{だいぶつ}}
\textbf{\ja{だいぶつ}} — \emph{daibutsu} — \textquotedblleft{}a great statue of the Buddha\textquotedblright{}.

\begin{quote}
\textbf{\ja{この} \ja{てらの} \ja{だいぶつは} \ja{たかいです。}} — \emph{kono tera no daibutsu wa takai desu} — \textquotedblleft{}The great statue at this temple is tall.\textquotedblright{}
\end{quote}

\begin{cousinweb}
\textbf{\ja{だいぶつ}} (\emph{daibutsu}) is Sino-Japanese (kango): two Chinese characters read with their on'yomi, readings borrowed from Chinese. They are \emph{dai}, \textquotedblleft{}large\textquotedblright{} (Middle Chinese \emph{dajH}) and \emph{butsu}, \textquotedblleft{}Buddha\textquotedblright{} (\emph{bjut}) (KANJIDIC2; Middle Chinese in Baxter's transcription).
\end{cousinweb}

\begin{grammarlens}[title={What you've built}]
One more word for getting around.
\end{grammarlens}

\subsection*{Your turn}
\begin{itemize}
\raggedright
  \item \emph{Say it:} \emph{daibutsu}
  \item \emph{Say it:} \emph{daibutsu}, once more
  \item \emph{Say it:} \emph{kono tera no daibutsu wa takai desu} — The great statue at this temple is tall.
  \item \emph{Recall:} say the Japanese for a visit to worship, then the Japanese for a statue of the Buddha, then say \emph{daibutsu} again
\end{itemize}

\begin{tcolorbox}[breakable,colback=teal!4,colframe=teal!35!black,title={Before you move on}]
What does \ja{だいぶつ} mean? (\textquotedblleft{}A great statue of the Buddha\textquotedblright{}.) Say it once more.
\end{tcolorbox}
\section[\texorpdfstring{\ja{ほんどう} (hondō)}{ほんどう (hondō)}]{\ja{ほんどう} (hondō) — the main hall of a temple}

\label{lesson:JA-C407-hondou}



\begin{quote}
Before the new one: say the Japanese for a statue of the Buddha, then the Japanese for a great statue of the Buddha.
\end{quote}

\subsection*{You'll want to know: \ja{ほんどう}}
\textbf{\ja{ほんどう}} — \emph{hondō} — \textquotedblleft{}the main hall of a temple\textquotedblright{}.

\begin{quote}
\textbf{\ja{ほんどうで} \ja{くつを} \ja{ぬぎました。}} — \emph{hondō de kutsu o nugimashita} — \textquotedblleft{}We took off our shoes in the main hall.\textquotedblright{}
\end{quote}

\begin{cousinweb}
\textbf{\ja{ほんどう}} (\emph{hondō}) is Sino-Japanese (kango): two Chinese characters read with their on'yomi, readings borrowed from Chinese. They are \emph{hon}, \textquotedblleft{}main\textquotedblright{} (Middle Chinese \emph{pwonX}) and \emph{dō}, \textquotedblleft{}hall\textquotedblright{} (\emph{dang}) (KANJIDIC2; Middle Chinese in Baxter's transcription).
\end{cousinweb}

\begin{grammarlens}[title={What you've built}]
One more word for getting around.
\end{grammarlens}

\subsection*{Your turn}
\begin{itemize}
\raggedright
  \item \emph{Say it:} \emph{hondō}
  \item \emph{Say it:} \emph{hondō}, once more
  \item \emph{Say it:} \emph{hondō de kutsu o nugimashita} — We took off our shoes in the main hall.
  \item \emph{Recall:} say the Japanese for a statue of the Buddha, then the Japanese for a great statue of the Buddha, then say \emph{hondō} again
\end{itemize}

\begin{tcolorbox}[breakable,colback=teal!4,colframe=teal!35!black,title={Before you move on}]
What does \ja{ほんどう} mean? (\textquotedblleft{}The main hall of a temple\textquotedblright{}.) Say it once more.
\end{tcolorbox}
\section[\texorpdfstring{\ja{おがむ} (ogamu)}{おがむ (ogamu)}]{\ja{おがむ} (ogamu) — to pray with joined hands}

\label{lesson:JA-C407-ogamu}



\begin{quote}
Before the new one: say the Japanese for a great statue of the Buddha, then the Japanese for the main hall of a temple.
\end{quote}

\subsection*{You'll want to know: \ja{おがむ}}
\textbf{\ja{おがむ}} — \emph{ogamu} — \textquotedblleft{}to pray with joined hands\textquotedblright{}.

\begin{quote}
\textbf{\ja{ほんどうの} \ja{まえで} \ja{おがみました。}} — \emph{hondō no mae de ogamimashita} — \textquotedblleft{}I put my hands together in front of the main hall.\textquotedblright{}
\end{quote}

Polite \textbf{\ja{おがみます}}, \emph{ogamimasu}; a godan verb. Past (polite): \textbf{\ja{そぼが} \ja{ほんどうで} \ja{おがみました。}}, \emph{sobo ga hondō de ogamimashita}.

\begin{cousinweb}
\textbf{\ja{おがむ}} (\emph{ogamu}) is a native Japanese word (wago), inherited rather than borrowed from Chinese; its kanji was attached later as a kun'yomi (JMdict).
\end{cousinweb}

\begin{grammarlens}[title={What you've built}]
One more word for getting around.
\end{grammarlens}

\subsection*{Your turn}
\begin{itemize}
\raggedright
  \item \emph{Say it:} \emph{ogamu}
  \item \emph{Say it:} \emph{ogamu}, once more
  \item \emph{Say it:} \emph{hondō no mae de ogamimashita} — I put my hands together in front of the main hall.
  \item \emph{Recall:} say the Japanese for a great statue of the Buddha, then the Japanese for the main hall of a temple, then say \emph{ogamu} again
\end{itemize}

\begin{tcolorbox}[breakable,colback=teal!4,colframe=teal!35!black,title={Before you move on}]
What does \ja{おがむ} mean? (\textquotedblleft{}To pray with joined hands\textquotedblright{}.) Say it once more.
\end{tcolorbox}
\section[(dialogue)]{Looking back: from A Day of Sightseeing to Inside the Temple}

\label{lesson:JA-R407-review-402-407}



\begin{quote}
Say the Japanese for tourism and a tourist.
\end{quote}

\subsection*{Your turn}
Say these aloud:

\begin{itemize}
\raggedright
  \item tourism, a tourist, a famous sight, a local speciality, to tour
  \item culture, a highlight, a study visit, a place visited on a trip, to visit a place
  \item scenery, to gaze at, to look out over, wonderful, to look down on
  \item a landscaped garden, famous, splendid, to become rough, to cool off
  \item a shrine gate, to pass under, shrine grounds, a Shinto priest, to offer up
  \item a visit to worship, a statue of the Buddha, a great statue of the Buddha, the main hall of a temple, to pray with joined hands
\end{itemize}

\begin{tcolorbox}[breakable,colback=teal!4,colframe=teal!35!black,title={Before you move on}]
Tourism? (\textbf{\ja{かんこう}}, \emph{kankō}.) A tourist? (\textbf{\ja{かんこうきゃく}}, \emph{kankōkyaku}.) A famous sight? (\textbf{\ja{めいしょ}}, \emph{meisho}.) A local speciality? (\textbf{\ja{めいぶつ}}, \emph{meibutsu}.) To tour? (\textbf{\ja{めぐる}}, \emph{meguru}.) Culture? (\textbf{\ja{ぶんか}}, \emph{bunka}.) A highlight? (\textbf{\ja{みどころ}}, \emph{midokoro}.) A study visit? (\textbf{\ja{けんがく}}, \emph{kengaku}.) A place visited on a trip? (\textbf{\ja{たびさき}}, \emph{tabisaki}.) To visit a place? (\textbf{\ja{おとずれる}}, \emph{otozureru}.) Scenery? (\textbf{\ja{ふうけい}}, \emph{fūkei}.) To gaze at? (\textbf{\ja{ながめる}}, \emph{nagameru}.) To look out over? (\textbf{\ja{みわたす}}, \emph{miwatasu}.) Wonderful? (\textbf{\ja{すばらしい}}, \emph{subarashii}.) To look down on? (\textbf{\ja{みおろす}}, \emph{miorosu}.) A landscaped garden? (\textbf{\ja{ていえん}}, \emph{teien}.) Famous? (\textbf{\ja{ゆうめい}}, \emph{yūmei}.) Splendid? (\textbf{\ja{みごと}}, \emph{migoto}.) To become rough? (\textbf{\ja{あれる}}, \emph{areru}.) To cool off? (\textbf{\ja{すずむ}}, \emph{suzumu}.) A shrine gate? (\textbf{\ja{とりい}}, \emph{torii}.) To pass under? (\textbf{\ja{くぐる}}, \emph{kuguru}.) Shrine grounds? (\textbf{\ja{けいだい}}, \emph{keidai}.) A Shinto priest? (\textbf{\ja{かんぬし}}, \emph{kannushi}.) To offer up? (\textbf{\ja{ささげる}}, \emph{sasageru}.) A visit to worship? (\textbf{\ja{さんぱい}}, \emph{sanpai}.) A statue of the Buddha? (\textbf{\ja{ぶつぞう}}, \emph{butsuzō}.) A great statue of the Buddha? (\textbf{\ja{だいぶつ}}, \emph{daibutsu}.) The main hall of a temple? (\textbf{\ja{ほんどう}}, \emph{hondō}.) To pray with joined hands? (\textbf{\ja{おがむ}}, \emph{ogamu}.)
\end{tcolorbox}
